Verifiable computation systems often need to prove large matrix multiplication
statements, but a direct SNARK arithmetization of a \(k \times k\) product
requires \(\mathcal{O}(k^3)\) constraints. Freivalds' randomized check reduces
the algebraic computation to vector-matrix products, but proving those products
inside a SNARK still costs \(\mathcal{O}(k^2)\) constraints.

We present \(\protocol\), a matrix-multiplication checking protocol that combines
Freivalds' randomized check with proximity testing over linear error-correcting
codes. The prover commits to encoded matrices and intermediate vectors before
the sampled query positions are derived. The CP-SNARK circuit then checks only the sampled codeword positions and commits to the values used inside the circuit, while Merkle openings and CP-Link proofs ensure consistency between the in-circuit witnesses and the externally committed values. We prove soundness for this committed-input setting under the soundness of the SNARK backend, the binding of the commitments, the correctness of the CP-Link checks, and \rev{the relative distance and unique-decoding radius of the Reed--Solomon code}.

For a fixed number \(t\) of sampled positions, the main in-circuit SNARK
relation has \revmath{\mathcal{O}(tk+E_{\mathsf{code}})} constraints, with additional
\revmath{\mathcal{O}(t\log n)+E_{\mathsf{link}}(k,t)} backend work for Merkle openings
and CP-Link checks. We implement \(\protocol\) in Go and compare it with a
Freivalds-based SNARK circuit. In the matrix benchmark at \(k=2^{12}\),
\(\protocol\) reduces the constraint count by \revmath{9.68\times} and shortens proof
generation time by \revmath{5.81\times}; verification \rev{remains approximately}
\revmath{0.13\,\text{s}} \rev{across the measured range}.
